\section{Introduction}\label{sec:intro}

Graph-based recommenders represent users and items as nodes and score a user--item
pair by propagating embeddings along the user--item graph. LightGCN
\citep{he2020lightgcn} is a standard backbone of this family. Such models are
usually trained on implicit feedback logs, in which an interaction is recorded only
when the item was shown and the user responded. The logged graph is therefore not
the preference graph: it records the exposure policy as much as the user's taste
\citep{marlin2009collaborative,chen2023bias}. A model fitted to such a log treats
every absence as informative, and the resulting bias is the subject of a large
literature on data missing not at random.

The pairwise loss that trains LightGCN, Bayesian personalised ranking (BPR)
\citep{rendle2009bpr}, draws its negatives uniformly from the items that the user
has no logged interaction with. Under exposure bias most of these items were never
shown to the user, and some of them would have been liked had they been shown.
These are false negatives, and they concentrate on rarely exposed items, where the
log carries the least information.
Matrix factorisation models have addressed this by treating exposure as a latent
variable \citep{liang2016exposure} or by re-weighting the likelihood with
propensities \citep{saito2020unbiased}. Loss-level debiasing with inverse
propensity or doubly robust (DR) estimators \citep{schnabel2016,wang2019drjl,guo2021mrdr}
is defined over observed or sampled entries. In our experiments a
pointwise DR loss on LightGCN (DR-LightGCN) reaches nDCG@5 of 0.5528 on Yahoo!\,R3,
well below BPR training of LightGCN on the same data (0.6708).

This paper starts from the two nuisance functions that a DR estimator already fits:
an exposure propensity $p_{ui}$ and an imputed preference $\hat Y_{ui}$. Under an independence
assumption (Section~\ref{sec:setting}), Bayes' rule applied to them gives the
probability that a pair absent from the log is a positive, conditional on its covariates,
\begin{equation}\label{eq:Q}
Q_{ui}=\Pr\bigl(Y_{ui}=1\mid O_{ui}Y_{ui}=0,\,x_{ui}\bigr)
=\frac{\hat Y_{ui}\,(1-p_{ui})}{1-p_{ui}\hat Y_{ui}}.
\end{equation}
The quantity $Q_{ui}$ is at most $\hat Y_{ui}$, and it is close to $\hat Y_{ui}$ when
the estimated exposure propensity is small. We use $Q$ in two places inside an otherwise unchanged LightGCN.
Each sampled BPR negative is down-weighted by $1-Q_{uj}$, so that likely false
negatives contribute less gradient. The standardised logit of the imputation is
added to the score as a fixed offset, whose weight is selected on validation data.
Both nuisance functions are fitted once on the biased log before training. Training
therefore uses no uniformly collected data, except the shipped Coat propensities, which the dataset authors estimated from a uniform sample (Section~\ref{sec:nuisance}). We call the resulting model \qp{}
(Q-weighted negatives with an imputation prior offset).

The evaluation design follows from a concern that progress claims in top-$N$
recommendation often do not survive a comparison with tuned simple baselines
\citep{dacrema2019progress,rendle2019baselines,anelli2023myth}. Some reported
debiasing gains rely on protocols that train on uniformly collected data. We use uniform data only for validation and testing (with the Coat exception noted in Section~\ref{sec:nuisance}). The
configuration of each split is selected on validation data, the per-user comparisons
are corrected with Holm's procedure \citep{holm1979}, and the selected configurations
are retrained with additional seeds. The results are mixed, and we report them as
such. \qp{} has a higher mean than every graph baseline on all four datasets (significant except
LightGCN-pt on KuaiRec), and it is the best method among those we ran on Yahoo!\,R3 and KuaiRand.
On Coat it is above iALS in the main run but not significantly different from it across
retraining seeds. On KuaiRec it is below DR-JL and MRDR in the main run and not significantly
different from them across seeds. On KuaiRec a
training-free propagation operator, DRUP, is the strongest method we ran.

DRUP, a training-free walk-corrected propagation operator, applies inverse-propensity and DR edge
estimates to multi-hop propagation and corrects the bias that repeated edges add to
walks of more than one hop. We keep DRUP as a baseline and as a component, and we
report a negative finding about its correction term. At matched configurations the
walk correction does not improve accuracy on the real datasets. It improves score
values only in a semi-synthetic setting with known propensities. Training-free DR
propagation remains the strongest method on KuaiRec, and the correction is not
responsible for that result: the same operator without the correction performs
equally well.

The contributions of the paper are as follows.
\begin{enumerate}
\item \emph{Method.} A closed-form posterior $Q_{ui}$ for pairs that are not logged
positives, computed from an inverse-propensity-weighted additive imputation and an
exposure propensity (Eq.~\eqref{eq:Q}). $Q$ is used twice in LightGCN: to weight
sampled BPR negatives and as a standardised imputation offset. The method adds one
selected hyper-parameter, the offset weight, and one lookup per sampled negative.
\item \emph{Evaluation.} A comparison on four public benchmarks with an unbiased test
set (Coat, Yahoo!\,R3, KuaiRand-Pure and KuaiRec), in which uniform data are used
only for validation and testing (Coat's shipped propensities excepted). The comparison uses per-user paired tests with Holm
correction within each dataset, and the main configurations are retrained with three
additional seeds.
\item \emph{Mechanism.} Ablations that replace $Q$ by the imputation alone, by item
popularity, or by a random permutation of its own values. The shuffled control
returns to the level of LightGCN, which shows that the weighting requires an
informative $Q$. The relative contribution of weighting and offset depends on the
dataset. On Coat, with unmatched grids (108 vs 54 configurations), the offset alone
recovers most of the gain.
\item \emph{A negative result on walk correction.} The idempotent correction of DRUP
does not improve accuracy on real data at matched configurations, and it improves
scores only under semi-synthetic conditions with known propensities.
\end{enumerate}

The rest of the paper is organised as follows. Section~\ref{sec:related} reviews the
related work and Table~\ref{tab:positioning} positions the method against the main
families. Section~\ref{sec:setting} defines the setting and notation. Section~\ref{sec:method}
describes the method and the DRUP baseline, Section~\ref{sec:setup} the datasets,
baselines and protocol, and Section~\ref{sec:results} the results. Section~\ref{sec:discussion}
discusses when the method helps and when it does not, together with its limitations,
and Section~\ref{sec:conclusion} concludes.
